% 本份微分几何独立笔记的写作约定（仅适用于本笔记）：
% 数学术语保留英文，其他叙述用中文；专业数学术语首次出现时附中文名称，之后不重复。
% 例如 Gaussian curvature（高斯曲率）；open set 等基础术语不用附中文名称。
% 基本保留用户原有措辞、表达和内容，不因片段不连贯而自行补写。
\documentclass[UTF8,11pt,a4paper]{ctexart}
\usepackage[margin=2.5cm]{geometry}
\usepackage{amsmath,amssymb}
\usepackage{hyperref}
\setlength{\parindent}{0pt}
\setlength{\parskip}{0.65em}
\title{Differential geometry of surfaces（曲面微分几何）}
\author{}
\date{}
\begin{document}
\maketitle
\section{The structure equations（结构方程）of $\mathbb{R}^n$}

设 $\{e_i\}_{i=1}^n$ 是 $\mathbb{R}^n$ 中或某个 open set $U$ 中的 $n$ 个 smooth vector fields（光滑向量场），并且构成一个 orthonormal moving frame（正交归一活动标架）。这就是说，
\[
\langle e_i,e_j\rangle=\delta_{ij}.
\]
考虑 $1$-form（一次微分形式）$\omega_j$，使得
\[
\omega_j(e_i)=\delta_{ij}.
\]
$\omega_j$ 是 smooth 的，这一点容易看出，因为可以写出 $\omega_j$ 关于 $dx_1,\ldots,dx_n$ 的分量表示。

另一方面，每个 vector field $e_i$ 都是一个 map
\[
e_i\colon U\longrightarrow\mathbb{R}^n.
\]
它在 $p\in U$ 处有 differential（微分）
\[
(de_i)_p\colon\mathbb{R}^n\longrightarrow\mathbb{R}^n.
\]
我们当然可以写
\[
(de_i)_p(v)=\sum_j(\omega_{ij})_p(v)e_j,
\]
这里作为分量的 $\omega_{ij}$ 是 smooth $1$-form。为简洁起见，写作
\[
de_i=\sum_j\omega_{ij}e_j.
\]
令 $e=(e_1,\ldots,e_n)^{\mathsf T}$，$W=(\omega_{ij})_{n\times n}$，则上式也可简写为 $de=We$。
称 $\omega_{ij}$ 为 \textbf{connection forms（联络形式）}。

这些 forms 并非全部 independent。例如，对
\[
\langle e_i,e_j\rangle=1
\]
取 differential，得到
\[
\langle de_i,e_j\rangle+\langle e_i,de_j\rangle=0,
\]
即
\[
\omega_{ij}+\omega_{ji}=0.
\]

对于 manifold（流形） $M$（以下经常是 $\mathbb{R}^3$ 中的 $2$ 维 surface（曲面）），如果有 inner product（内积）
\[
\langle v,u\rangle_p,\qquad v,u\in T_pM,
\]
也可定义 orthonormal moving frame 等。

这里需要区分 $d$ 的两种用法。最初 $e_i\colon U\to\mathbb{R}^n$ 是 map，其 differential 在每点为
\[
(de_i)_p\colon T_pU\longrightarrow\mathbb{R}^n.
\]
严格来说，map 的 differential 的 target 是 $T_{e_i(p)}\mathbb{R}^n$；这里通过 Euclidean space 的自然识别，将它统一视为固定的 $\mathbb{R}^n$。因此，不同 $p$ 处的 $(de_i)_p$ 的像都落在同一个 $\mathbb{R}^n$ 中，而不是保留在随点变化的 tangent spaces 中。

将这些 linear maps 合在一起，$de_i$ 是一个 $\mathbb{R}^n$-valued $1$-form（向量值一次微分形式）。对它再写 $d(de_i)$ 时，$d$ 指的是作用于 vector-valued forms 的 exterior derivative，而不是把 $de_i$ 当作普通 map 再求 differential。该 exterior derivative 可逐分量定义，并满足 $d^2=0$；它在 $\mathbb{R}^n$-valued functions 上与原来的 differential 一致。因此
\[
d(de_i)=0.
\]
换成 moving frame 表示 $de_i=\sum_j\omega_{ij}e_j$ 后，$e_j$ 也随点变化，求 exterior derivative 时必须一并求导：
\[
d(de_i)=\sum_j(d\omega_{ij})e_j-\sum_j\omega_{ij}\wedge de_j.
\]

\textbf{Theorem（structure equations）}
\[
d\omega_i=\sum_k\omega_k\wedge\omega_{ki},
\qquad
d\omega_{ij}=\sum_k\omega_{ik}\wedge\omega_{kj}.
\]

\textbf{Proof.} 令 $x\colon U\longrightarrow\mathbb{R}^n$ 是 inclusion（包含映射）。
\[
dx=\sum_i\omega_i e_i,
\qquad
de_i=\sum_j\omega_{ij}e_j.
\]
于是
\begin{align*}
d^2x
&=\sum_i(d\omega_i)e_i-\sum_i\omega_i\wedge de_i\\
&=\sum_i(d\omega_i)e_i
  -\sum_i\omega_i\wedge\left(\sum_j\omega_{ij}e_j\right)=0.
\end{align*}
比较 $e_i$ 的系数，得
\[
d\omega_i=\sum_k\omega_k\wedge\omega_{ki}.
\]
又有
\begin{align*}
d^2e_i
&=d\left(\sum_j\omega_{ij}e_j\right)\\
&=\sum_j d\omega_{ij}e_j
  -\sum_j\omega_{ij}\wedge\left(\sum_k\omega_{jk}e_k\right)=0.
\end{align*}
同理得到
\[
d\omega_{ij}=\sum_k\omega_{ik}\wedge\omega_{kj}.
\]

\medskip
现在假设 $M^n$ 是 $\mathbb{R}^{n+k}$ 的 submanifold（子流形）。

以后 frame 都指之前的 orthonormal moving frame。我们称一个 frame 是 adapted（适配）的，如果 restrict 到 $M^n$ 的 tangent space（切空间） 里，那 $e_1,\ldots,e_n$ 就是 $M^n$ 的 frame。设 inclusion map 是 $x$，则有 $M$ 上的
\[
x^*(\omega_i),\qquad x^*(\omega_{ij}).
\]
另一方面，由 submanifold 的定义，对于 $M$ 上的 frame，总能（至少是 locally）扩展成 $\mathbb{R}^{n+k}$ 中的 frame。所以先前的讨论方法也是最一般的。

\medskip
\textbf{Lemma.} 设 $V$ 是 $\mathbb{R}$ 上的 vector space，$\dim V=n$，$\omega_1,\ldots,\omega_r$（$r\leq n$）是 linearly independent 的 $1$-forms。如果 forms
\[
\theta_1,\ldots,\theta_r\colon V\longrightarrow\mathbb{R}
\]
满足
\[
\sum_{i=1}^r\omega_i\wedge\theta_i=0,
\]
则
\[
\theta_i=\sum_{j=1}^r a_{ij}\omega_j,
\qquad a_{ij}=a_{ji}.
\]

\textbf{Proof.} 设 $\omega_1,\ldots,\omega_r$ 与 $\omega_{r+1},\ldots,\omega_n$ 组成 $V^*$ 的 basis。写作
\[
\theta_i=\sum_{j=1}^r a_{ij}\omega_j
+\sum_{\ell=r+1}^n b_{i\ell}\omega_\ell.
\]
现在代入条件，得到
\[
\sum_{i=1}^r\omega_i\wedge
\left[
\left(\sum_{p=1}^r a_{ip}\omega_p\right)
+\sum_{q=r+1}^n b_{iq}\omega_q
\right]=0.
\]
比较系数即可（注意 $\omega_s\wedge\omega_\ell$，$s<\ell$，是 $\Lambda^2(V)$ 的 basis）。

\medskip
\textbf{Lemma.} 设 $U\subset\mathbb{R}^n$，$\omega_1,\ldots,\omega_n$ 是 $U$ 上 linearly independent 的 $1$-forms。如果存在 $1$-forms $\omega_{ij}$（$i,j=1,\ldots,n$），满足
\[
\omega_{ij}+\omega_{ji}=0,
\qquad d\omega_j=\sum_k\omega_k\wedge\omega_{kj},
\]
则对于所有 $k,j$，$\omega_{kj}$ 唯一。

可以想到，之后这个 lemma 会被用来证明一些 $M$ 上 $\omega_{ij}$ 的 uniqueness（唯一性）。

\textbf{Proof.} 如果另有 $\overline{\omega}_{ij}$ 满足
\[
\overline{\omega}_{ij}+\overline{\omega}_{ji}=0,
\qquad d\omega_j=\sum_k\omega_k\wedge\overline{\omega}_{kj},
\]
则
\[
0=d\omega_j-d\omega_j
=\sum_k\omega_k\wedge(\omega_{kj}-\overline{\omega}_{kj}).
\]
由上一个 lemma，对于 $k=1,\ldots,n$，有
\[
\omega_{kj}-\overline{\omega}_{kj}=\sum_r B^j_{kr}\omega_r,
\]
其中
\begin{equation}
B^j_{kr}=B^j_{rk}.
\tag{S}\label{eq:dg-B-symmetry}
\end{equation}
现在再看 $\omega_{kj}+\omega_{jk}=0$ 的条件：
\[
0=\omega_{kj}+\omega_{jk}
-\overline{\omega}_{kj}-\overline{\omega}_{jk}
=\sum_r(B^j_{kr}+B^k_{jr})\omega_r.
\]
由 $\omega_r$ 的 linear independence，对所有 $r$ 有
\begin{equation}
B^j_{kr}=-B^k_{jr}.
\tag{A}\label{eq:dg-B-antisymmetry}
\end{equation}
因此，最后的等式链可写成
\begin{align*}
B^j_{kr}
&\overset{\text{(S)}}{=}B^j_{rk}
 \overset{\text{(A)}}{=}-B^r_{jk}
 \overset{\text{(S)}}{=}-B^r_{kj}\\
&\overset{\text{(A)}}{=}B^k_{rj}
 \overset{\text{(S)}}{=}B^k_{jr}
 \overset{\text{(A)}}{=}-B^j_{kr}.
\end{align*}
这里每一步均将 (S) 或 (A) 中的指标相应替换；(S) 交换两个下标，(A) 交换上标与第一个下标并变号。
故 $B^j_{kr}=0$，从而 $\omega_{kj}=\overline{\omega}_{kj}$。

\section{$\mathbb{R}^3$ 中的 surfaces}
设 surface $M\subset\mathbb{R}^3$。在 $\mathbb{R}^3$ 的一个 open set 中选取 frame $e_1,e_2,e_3$，相应的 forms 为
\[
\omega_1,\omega_2,\omega_3,\qquad
\omega_{12},\omega_{13},\omega_{23}.
\]
由于 $\omega_{ij}+\omega_{ji}=0$，有 $\omega_{ii}=0$。现在写出已知的 structure equations：
\begin{align*}
d\omega_1&=\omega_2\wedge\omega_{21}+\omega_3\wedge\omega_{31},\\
d\omega_2&=\omega_1\wedge\omega_{12}+\omega_3\wedge\omega_{32},\\
d\omega_3&=\omega_1\wedge\omega_{13}+\omega_2\wedge\omega_{23},\\
d\omega_{12}&=\omega_{13}\wedge\omega_{32},\\
d\omega_{13}&=\omega_{12}\wedge\omega_{23},\\
d\omega_{23}&=\omega_{21}\wedge\omega_{13}.
\end{align*}

回忆：由 inclusion $x\colon M\hookrightarrow\mathbb{R}^3$ 诱导的 pullback（拉回） $x^*$ 与 exterior derivative（外微分）$d$ 和 wedge product（楔积）可交换，即
\[
x^*(d\alpha)=d(x^*\alpha),\qquad
x^*(\alpha\wedge\beta)=x^*\alpha\wedge x^*\beta.
\]
这里选取 adapted frame，使 $e_1,e_2$ 是 $M$ 上的 tangent frame（切标架），$e_3$ 是 unit normal（单位法向量）。于是
\[
x^*(\omega_3)=0,
\]
因为 $T_pM\perp e_3(p)$。

为简洁起见，下面仍用 $\omega_i,\omega_{ij}$ 表示各自的 pullback。由上面的 structure equations 得到 $M$ 上的方程：
\begin{align*}
d\omega_1&=\omega_2\wedge\omega_{21}
          =-\omega_2\wedge\omega_{12},\\
d\omega_2&=\omega_1\wedge\omega_{12},\\
0=d\omega_3&=\omega_1\wedge\omega_{13}+\omega_2\wedge\omega_{23},\\
d\omega_{12}&=\omega_{13}\wedge\omega_{32}
             =-\omega_{13}\wedge\omega_{23},\\
d\omega_{13}&=\omega_{12}\wedge\omega_{23},\\
d\omega_{23}&=\omega_{21}\wedge\omega_{13}
             =-\omega_{12}\wedge\omega_{13}.
\end{align*}

由
\[
\omega_1\wedge\omega_{13}+\omega_2\wedge\omega_{23}=0,
\]
使用上一个关于 $1$-forms 的 lemma，得到
\[
\omega_{13}=h_{11}\omega_1+h_{12}\omega_2,
\qquad
\omega_{23}=h_{21}\omega_1+h_{22}\omega_2,
\qquad h_{12}=h_{21}.
\]

\subsection*{Structure equations 的进一步化简}
记 Gaussian curvature（高斯曲率）与 mean curvature（平均曲率） 分别为
\[
K=h_{11}h_{22}-h_{12}^2,
\qquad H=\frac12(h_{11}+h_{22}).
\]
于是
\begin{align*}
d\omega_{12}
&=\omega_{13}\wedge\omega_{32}\\
&=(h_{11}\omega_1+h_{12}\omega_2)
\wedge(-h_{12}\omega_1-h_{22}\omega_2)\\
&=(-h_{11}h_{22}+h_{12}^2)\omega_1\wedge\omega_2
=-K\omega_1\wedge\omega_2.
\end{align*}
另外，
\begin{align*}
\omega_{13}\wedge\omega_2+\omega_1\wedge\omega_{23}
&=(h_{11}\omega_1+h_{12}\omega_2)\wedge\omega_2\\
&\quad+\omega_1\wedge(h_{21}\omega_1+h_{22}\omega_2)\\
&=2H\omega_1\wedge\omega_2.
\end{align*}

\subsection*{$h_{ij}$ 的解释}
考虑 Gauss map（高斯映射） $e_3\colon U\longrightarrow S^2$，其中 $U$ 是 $M$ 上选定 unit normal 的局部区域。有
\begin{align*}
de_3
&=\sum_i\omega_{3i}e_i
=\omega_{31}e_1+\omega_{32}e_2\\
&=-(h_{11}\omega_1+h_{12}\omega_2)e_1
  -(h_{21}\omega_1+h_{22}\omega_2)e_2.
\end{align*}
现在对 $v=a_1e_1(p)+a_2e_2(p)\in T_pM$ 代入，得到
\[
(de_3)_p(v)
=-(h_{11}(p)a_1+h_{12}(p)a_2)e_1(p)
 -(h_{21}(p)a_1+h_{22}(p)a_2)e_2(p).
\]
将 $T_{e_3(p)}S^2=e_3(p)^\perp$ 与 $T_pM$ 自然识别后，在 basis $(e_1(p),e_2(p))$ 下，
\[
\bigl[(de_3)_p(v)\bigr]
=-\begin{pmatrix}
h_{11}(p)&h_{12}(p)\\
h_{21}(p)&h_{22}(p)
\end{pmatrix}
\begin{pmatrix}a_1\\a_2\end{pmatrix}.
\]

\textbf{Corollary.} $(de_3)_p$ 是 self-adjoint operator（自伴算子），因此对于 $v,u\in T_pM$，
\[
\langle(de_3)_p(v),u\rangle
=\langle v,(de_3)_p(u)\rangle.
\]
而且 $(de_3)_p$ 有两个 real eigenvalues（实特征值）（计 multiplicity（重数）），设为 $-\lambda_1,-\lambda_2$。

因此，在上述 tangent spaces 的自然识别下，
\[
K(p)=\det((de_3)_p)=\lambda_1\lambda_2,
\qquad
H(p)=-\frac12\operatorname{tr}((de_3)_p)
=\frac{\lambda_1+\lambda_2}{2}.
\]
\medskip
\textbf{First fundamental form（第一基本形式）.} 对 adapted frame $\{e_i\}$，$i=1,2,3$，
\[
\mathrm{I}(v)=\langle v,v\rangle
=(\omega_1\omega_1+\omega_2\omega_2)(v).
\]
这里 $\omega_1\omega_1$ 就是 $\omega_1(v)\cdot\omega_1(v)$。

\textbf{Second fundamental form（第二基本形式）.}
\begin{align*}
\mathrm{II}_p(v)
&=-\langle de_3(v),v\rangle_p\\
&=-\langle\omega_{31}(v)e_1+\omega_{32}(v)e_2,v\rangle_p\\
&=-(\omega_{31}\omega_1+\omega_{32}\omega_2)(v)\\
&=(\omega_{13}\omega_1+\omega_{23}\omega_2)(v).
\end{align*}
（回忆：$v=\omega_1(v)e_1+\omega_2(v)e_2+\omega_3(v)e_3$，这里 $\omega_3(v)=0$。）

\subsection*{另一种解释}
$\alpha\colon(-\varepsilon,\varepsilon)\longrightarrow M$ 是 arc length curve（弧长参数曲线），以 $s$ 参数化，$\alpha(0)=p$，$\alpha'(0)=v\in T_pM$。当然有 $\alpha'(s)\perp e_3(s)$。我们将看到 $\mathrm{II}_p(v)$ 是 $\dfrac{d^2\alpha}{ds^2}$ 在 normal direction（法向）上的投影：
\begin{align*}
\left.\left\langle\frac{d^2\alpha}{ds^2},e_3(s)\right\rangle\right|_{s=0}
&=-\left.\left\langle\frac{d\alpha}{ds},\frac{de_3}{ds}\right\rangle\right|_{s=0}\\
&=-\langle\omega_1e_1+\omega_2e_2,
\omega_{31}e_1+\omega_{32}e_2\rangle(v)
=\mathrm{II}_p(v).
\end{align*}
称 $p$ 处 $v$ 方向的 normal curvature（法曲率）为
\[
k_n(v)=\mathrm{II}_p(v).
\]

由于 $de_3$ 是 self-adjoint 的，设 $\|v\|=1$，$v\in T_pM$，
\[
v=\cos\theta\,u_1+\sin\theta\,u_2,
\]
其中 $u_1,u_2$ 是 $de_3$ 的两个 orthonormal eigenvectors（正交归一特征向量），$de_3(u_i)=-\lambda_i u_i$；对应的 eigenvalues $-\lambda_1,-\lambda_2$ 可能相等。从而
\[
k_n(v)=\lambda_1\cos^2\theta+\lambda_2\sin^2\theta.
\]
\textbf{Corollary.} $\mathrm{II}_p(v)$ 在 $\|v\|=1$ 上的两个极值是 $\lambda_1,\lambda_2$。

称 $\lambda_i$ 为 principal curvatures（主曲率），$u_i$ 为 principal directions（主方向）。


\medskip
\textbf{Theorem.} 设 $U,U'$ 是 $\mathbb{R}^3$ 中两个 connected 的 $2$ 维 submanifolds。假设 $U$ 上存在 adapted frame $\{e_i\}$，$U'$ 上存在 adapted frame $\{e_i'\}$，$i=1,2,3$，以及 diffeomorphism $f\colon U\to U'$，使得
\[
f^*\omega_i'=\omega_i,
\qquad
f^*\omega_{ij}'=\omega_{ij},
\qquad i,j=1,2,3.
\]
则存在 rigid motion（刚体运动）$\rho\colon\mathbb{R}^3\to\mathbb{R}^3$，使得 $\rho|_U=f$。
\section{曲面的 intrinsic geometry（内蕴几何）}
以下我们总设 $M$ 是一个 $2$ 维 Riemannian manifold（黎曼流形）（直观上就是已经有了光滑的 $\mathrm{I}_p(v)$，但是已经看不到 $e_3$）。

这时候，对于 frame $\{e_i\}$、coframe $\{\omega_i\}$，存在 $1$-forms $\omega_{12},\omega_{21}$，使得
\[
\omega_{12}+\omega_{21}=0,
\qquad
d\omega_1=\omega_{12}\wedge\omega_2,
\qquad
d\omega_2=\omega_{21}\wedge\omega_1.
\]
（注意这里就没有良定义的 $de_i$ 了。）

\textbf{Proof.} 定义
\[
\omega_{12}(e_1)=d\omega_1(e_1,e_2),
\qquad
\omega_{12}(e_2)=d\omega_2(e_1,e_2)
\]
即可。

设 $\{\bar e_1,\bar e_2\}$ 是 $U$ 中的另一个 frame。如果 $\{\bar e_1,\bar e_2\}$ 与 $\{e_1,e_2\}$ 有相同的 orientation（定向），则
\[
\bar e_1=\cos\theta\,e_1+\sin\theta\,e_2,
\qquad
\bar e_2=-\sin\theta\,e_1+\cos\theta\,e_2,
\]


相反则是
\[
\bar e_1=\cos\theta\,e_1+\sin\theta\,e_2,
\qquad
\bar e_2=\sin\theta\,e_1-\cos\theta\,e_2,
\]
其中 $\theta$ 是局部选取的 smooth function，必要时缩小 $U$。

\textbf{Proposition.} 设 $\{e_1,e_2\}$ 与 $\{\bar e_1,\bar e_2\}$ 是 $U$ 上的两个 frames，$\omega_{12},\bar\omega_{12}$ 为各自的 connection forms。按上述 frame 变换中的 $\theta$，定义
\[
\tau=d\theta.
\]
若两个 frames 同定向，则
\[
\omega_{12}=\bar\omega_{12}-\tau;
\]
若反定向，则
\[
\omega_{12}=-\bar\omega_{12}-\tau.
\]
\textbf{Proof（同定向情形）.} 记 $c=\cos\theta$、$s=\sin\theta$。由 frame 的变换，
\[
\omega_1=c\bar\omega_1-s\bar\omega_2,
\qquad
\omega_2=s\bar\omega_1+c\bar\omega_2.
\]
利用 $\bar\omega_{21}=-\bar\omega_{12}$ 及带 bar 的 structure equations，分组得到
\begin{align*}
d\omega_1
&=-s\,d\theta\wedge\bar\omega_1
+c\,\bar\omega_{12}\wedge\bar\omega_2\\
&\quad-c\,d\theta\wedge\bar\omega_2
-s\,\bar\omega_{21}\wedge\bar\omega_1\\
&=(\bar\omega_{12}-d\theta)
\wedge(s\bar\omega_1+c\bar\omega_2)\\
&=(\bar\omega_{12}-d\theta)\wedge\omega_2.
\end{align*}
同样，
\begin{align*}
d\omega_2
&=c\,d\theta\wedge\bar\omega_1
+s\,\bar\omega_{12}\wedge\bar\omega_2\\
&\quad-s\,d\theta\wedge\bar\omega_2
+c\,\bar\omega_{21}\wedge\bar\omega_1\\
&=-(\bar\omega_{12}-d\theta)
\wedge(c\bar\omega_1-s\bar\omega_2)\\
&=-(\bar\omega_{12}-d\theta)\wedge\omega_1.
\end{align*}
因此，令 $\eta_{12}=\bar\omega_{12}-d\theta$、$\eta_{21}=-\eta_{12}$，就有
\[
\eta_{12}+\eta_{21}=0,
\qquad
d\omega_1=\eta_{12}\wedge\omega_2,
\qquad
d\omega_2=\eta_{21}\wedge\omega_1.
\]
这与 $\omega_{12},\omega_{21}$ 对同一个 coframe 满足的 structure equations 相同。由之前的 \textbf{connection forms uniqueness lemma}，
\[
\omega_{12}=\eta_{12}=\bar\omega_{12}-d\theta
=\bar\omega_{12}-\tau.
\]

\textbf{Theorem.} 设 $\{e_1,e_2\}$ 与 $\{\bar e_1,\bar e_2\}$ 是 $U$ 上同定向的两个 frames，局部地以之前的 angle function $\theta$ 表示它们的变换。设 $\gamma\colon I\to U$ 是 curve，$\gamma(t_0)=p$，并选定初始有向角 $\varphi_0$，使
\[
\bar e_1(p)=\cos\varphi_0\,e_1(p)+\sin\varphi_0\,e_2(p).
\]
则整个参数区间 $I$ 上存在唯一的 differentiable angle function $\varphi$，满足 $\varphi(t_0)=\varphi_0$ 及
\[
\bar e_1(\gamma(t))
=\cos\varphi(t)\,e_1(\gamma(t))
+\sin\varphi(t)\,e_2(\gamma(t)).
\]
其 differential 满足
\[
d\varphi=\gamma^*(d\theta),
\qquad
\varphi(t)=\varphi_0+\int_{t_0}^{t}
(d\theta)_{\gamma(s)}\bigl(\gamma'(s)\bigr)\,ds.
\]
这里 $\theta$ 只需局部选取；不同局部选择在重叠处相差局部常数 $2\pi m$，因此 $d\theta$ 一致，定义了 $U$ 上的同一个 $1$-form。$\varphi$ 是沿 curve 的累计有向转角，不必限制在某个 $2\pi$ 长的范围内。

\textbf{Proposition.} 设 $M$ 是一个 $2$ 维 Riemannian manifold。对于每个 $p\in M$，在 $p$ 附近选择 moving frame $\{e_1,e_2\}$，通过
\[
(d\omega_{12})_p=-K(p)(\omega_1\wedge\omega_2)_p
\]
定义一个数 $K(p)$。则 $K(p)$ 不依赖于 frame 的选择，称为 $M$ 在 $p$ 处的 Gaussian curvature。

\textbf{Proof.} 设 $\{\bar e_1,\bar e_2\}$ 是 $p$ 附近的另一个 moving frame。先假设两个 frames 同定向。由前面的 Proposition，局部地有
\[
\omega_{12}=\bar\omega_{12}-d\theta.
\]
由于 $d(d\theta)=0$，有 $d\omega_{12}=d\bar\omega_{12}$。又由 coframe 的变换，
\[
\bar\omega_1\wedge\bar\omega_2=\omega_1\wedge\omega_2.
\]
从而
\begin{align*}
-K\omega_1\wedge\omega_2
&=d\omega_{12}=d\bar\omega_{12}\\
&=-\bar K\bar\omega_1\wedge\bar\omega_2
=-\bar K\omega_1\wedge\omega_2.
\end{align*}
由于 $\omega_1\wedge\omega_2$ 处处非零，得到 $K=\bar K$。

若两个 frames 反定向，则
\[
\omega_{12}=-\bar\omega_{12}-d\theta,
\qquad
\bar\omega_1\wedge\bar\omega_2=-\omega_1\wedge\omega_2.
\]
因此
\[
-K\omega_1\wedge\omega_2
=d\omega_{12}=-d\bar\omega_{12}
=\bar K\bar\omega_1\wedge\bar\omega_2
=-\bar K\omega_1\wedge\omega_2,
\]
同样得到 $K=\bar K$。\hfill$\square$

\medskip
\textbf{Definition.} 设 $M$ 是一个 $2$ 维 Riemannian manifold，$Y$ 是 $M$ 上的 differentiable vector field。取 $p\in M$、$x\in T_pM$，并考虑 curve $\alpha\colon(-\varepsilon,\varepsilon)\to M$，满足 $\alpha(0)=p$、$\alpha'(0)=x$。为定义 $Y$ 在 $p$ 处沿 $x$ 的 covariant derivative（协变导数）$(\nabla_xY)(p)$，在 $p$ 附近选 moving frame $\{e_i\}$，将 $Y(\alpha(t))$ 在这个 frame 下表示为
\[
Y(\alpha(t))=\sum_i y_i(t)e_i,\qquad i=1,2,
\]
并令
\[
(\nabla_xY)(p)
=\sum_i\left(
\frac{dy_i}{dt}(0)+\sum_j\omega_{ji}(x)y_j(0)
\right)e_i,
\qquad i,j=1,2,
\]
其中约定 $\omega_{ii}=0$。

在二维情形，展开为
\[
(\nabla_xY)(p)=\bigl(y_1'(0)-\omega_{12}(x)y_2(0)\bigr)e_1(p)+\bigl(y_2'(0)+\omega_{12}(x)y_1(0)\bigr)e_2(p).
\]
其中 prime 表示对 curve 的参数 $t$ 求导。

\par\noindent

$\omega_{12}$ 与 covariant derivative 的关系为
\[
\omega_{12}(x)
=\langle\nabla_xe_1,e_2\rangle
=-\langle e_1,\nabla_xe_2\rangle.
\]
事实上，$\nabla_xe_1=\omega_{12}(x)e_2$，$\nabla_xe_2=-\omega_{12}(x)e_1$。

\par\noindent\textbf{Lemma.} Covariant derivative 不依赖于 frame 的选择。

\textbf{Proof（同定向情形）.} 沿 $\alpha(t)$ 计算，以下各量最终均在 $t=0$ 取值。记
\[
c=\cos\theta,\qquad s=\sin\theta,\qquad
\theta'=(\theta\circ\alpha)'(0)=d\theta_p(x),\qquad
A=\bar\omega_{12}(x).
\]
由同定向的 frame 变换及 $Y=y_1e_1+y_2e_2=\bar y_1\bar e_1+\bar y_2\bar e_2$，有
\[
\begin{aligned}
e_1&=c\bar e_1-s\bar e_2,& e_2&=s\bar e_1+c\bar e_2,\\
y_1&=c\bar y_1-s\bar y_2,& y_2&=s\bar y_1+c\bar y_2.
\end{aligned}
\]
对分量求导，得到
\begin{align*}
y_1'&=-s\theta'\bar y_1+c\bar y_1'-c\theta'\bar y_2-s\bar y_2',\\
y_2'&=c\theta'\bar y_1+s\bar y_1'-s\theta'\bar y_2+c\bar y_2'.
\end{align*}
由之前的 connection form 变换公式，$\omega_{12}(x)=A-\theta'$。将这些表达式代入
\[
\nabla_xY=(y_1'-\omega_{12}(x)y_2)e_1
+(y_2'+\omega_{12}(x)y_1)e_2.
\]
换成 $\bar e_1,\bar e_2$ 后，$\bar e_1$ 的系数为
\[
C_1=c\bigl(y_1'-(A-\theta')y_2\bigr)
+s\bigl(y_2'+(A-\theta')y_1\bigr).
\]
全部代入，尚未化简时为
\begin{align*}
C_1={}&c\bigl[-s\theta'\bar y_1+c\bar y_1'
-c\theta'\bar y_2-s\bar y_2'
-(A-\theta')(s\bar y_1+c\bar y_2)\bigr]\\
&+s\bigl[c\theta'\bar y_1+s\bar y_1'
-s\theta'\bar y_2+c\bar y_2'
+(A-\theta')(c\bar y_1-s\bar y_2)\bigr].
\end{align*}
按 $\bar y_1',\bar y_2',\bar y_1,\bar y_2$ 收集系数，分别为
\begin{align*}
\bar y_1':&\quad c^2+s^2=1,\\
\bar y_2':&\quad -cs+sc=0,\\
\bar y_1:&\quad -cs\theta'-cs(A-\theta')
+sc\theta'+sc(A-\theta')=0,\\
\bar y_2:&\quad -c^2\theta'-c^2(A-\theta')
-s^2\theta'-s^2(A-\theta')=-A.
\end{align*}
因此 $C_1=\bar y_1'-A\bar y_2$。同样，$\bar e_2$ 的系数为
\begin{align*}
C_2&=-s\bigl(y_1'-(A-\theta')y_2\bigr)
+c\bigl(y_2'+(A-\theta')y_1\bigr)\\
&=(-sy_1'+cy_2')+(A-\theta')(sy_2+cy_1)\\
&=(\bar y_2'+\theta'\bar y_1)+(A-\theta')\bar y_1
=\bar y_2'+A\bar y_1.
\end{align*}
故
\[
\nabla_xY
=\bigl(\bar y_1'-\bar\omega_{12}(x)\bar y_2\bigr)\bar e_1
+\bigl(\bar y_2'+\bar\omega_{12}(x)\bar y_1\bigr)\bar e_2,
\]
与用带 bar 的 frame 定义的结果一致。同定向情形得证。


\medskip
在进入之后的概念前，先看一下 $\mathbb{R}^2$ 作为 $M$ 有什么把握。
\begin{enumerate}
\item $K\equiv0$。注意由之前的论述，$K$ 无关 $e_i$ 的选择。令 $e_1=\partial_x$、$e_2=\partial_y$，从而 $\omega_1=dx$，于是 $d\omega_1=0$、$d\omega_2=0$。由满足
\[
d\omega_1=\omega_{12}\wedge\omega_2,
\qquad d\omega_2=\omega_{21}\wedge\omega_1
\]
的 $\omega_{12}$ 的唯一性，得到 $\omega_{12}=0$，从而 $K=0$。

\item $\nabla_xY$。仍假设 $e_1=\partial_x$、$e_2=\partial_y$，那 $\omega_{12}=0$，就得
\[
\nabla_xY=y_1'e_1+y_2'e_2,
\]
即方向导数。
\end{enumerate}

\medskip
\textbf{Proposition.} 设 $V$ 是沿 curve $\alpha(s)$ 的 differentiable tangent vector field，则
\[
\frac{d}{ds}\langle V,V\rangle
=\langle\nabla_{\alpha'}V,V\rangle
+\langle V,\nabla_{\alpha'}V\rangle.
\]

\textbf{Proof.} 沿 $\alpha$ 写 $V=v_1e_1+v_2e_2$，记 $a=\omega_{12}(\alpha')$。由于 frame 是 orthonormal，
\[
\frac{d}{ds}\langle V,V\rangle
=\frac{d}{ds}(v_1^2+v_2^2)
=2v_1v_1'+2v_2v_2'.
\]
另一方面，
\[
\nabla_{\alpha'}V=(v_1'-av_2)e_1+(v_2'+av_1)e_2,
\]
所以
\begin{align*}
\langle\nabla_{\alpha'}V,V\rangle
&=(v_1'-av_2)v_1+(v_2'+av_1)v_2\\
&=v_1v_1'+v_2v_2'.
\end{align*}
两个 connection 项由 $\omega_{21}=-\omega_{12}$ 恰好抵消。再由 inner product 的对称性，结论成立。\hfill$\square$
\textbf{Definition.} 沿 curve $\alpha\colon I\to M$ 的 vector field $Y$ 称为沿 $\alpha$ 的 parallel vector field（平行向量场），如果
\[
\nabla_{\alpha'(t)}Y=0,\qquad t\in I.
\]
由上面的 Proposition，若 $V$ 沿 $\alpha$ 平行，则 $\dfrac{d}{ds}\|V\|^2=0$，因此 $\|V\|$ 恒定。


\textbf{Proposition.} 设 $M\subset\mathbb{R}^3$ 采用 Euclidean inner product 诱导的 metric，$Y$ 是沿 $\alpha\colon I\to M$ 的 tangent vector field。则
\[
\nabla_{\alpha'}Y=\left(\frac{dY}{dt}\right)^{\top},
\]
其中 ${}^{\top}$ 表示向 $T_{\alpha(t)}M$ 的 orthogonal projection（正交投影）。因此，$Y$ 沿 $\alpha$ 平行，当且仅当 $dY/dt$ 的 tangent projection 为零，即 $dY/dt$ 沿 surface normal（包括零向量的情形）。

\textbf{Proof.} 沿 $\alpha$ 局部选取 adapted frame $e_1,e_2,e_3$，$e_3$ 为 unit normal。
沿 curve $\alpha(t)$ 时，$e_i(t)$ 实际表示 $e_i\circ\alpha(t)$。由 chain rule（链式法则），
\[
\frac{de_i}{dt}
=(de_i)_{\alpha(t)}\bigl(\alpha'(t)\bigr).
\]
因此，将 $de_i=\sum_j\omega_{ij}e_j$ 作用到 $\alpha'(t)$ 上，就得到
\[
\frac{de_i}{dt}=\sum_j\omega_{ij}(\alpha')e_j.
\]
这里 $\omega_{ij}$ 是 $1$-form，而 $\omega_{ij}(\alpha')$ 已经是 scalar function。在 $\mathbb{R}^3$ 中，具体为
\[
\frac{de_1}{dt}=\omega_{12}(\alpha')e_2+\omega_{13}(\alpha')e_3,
\qquad
\frac{de_2}{dt}=-\omega_{12}(\alpha')e_1+\omega_{23}(\alpha')e_3.
\]
所以，沿 surface 上的 curve 写 $Y=y_1e_1+y_2e_2$ 时，普通的乘积求导公式给出
\begin{align*}
\frac{dY}{dt}
&=y_1'e_1+y_2'e_2+y_1\frac{de_1}{dt}+y_2\frac{de_2}{dt}\\
&=y_1'e_1+y_2'e_2
+y_1\bigl[\omega_{12}(\alpha')e_2+\omega_{13}(\alpha')e_3\bigr]\\
&\quad+y_2\bigl[-\omega_{12}(\alpha')e_1+\omega_{23}(\alpha')e_3\bigr].
\end{align*}


按 $e_1,e_2,e_3$ 收集系数，得到
\begin{align*}
\frac{dY}{dt}
&=\bigl(y_1'-\omega_{12}(\alpha')y_2\bigr)e_1
+\bigl(y_2'+\omega_{12}(\alpha')y_1\bigr)e_2\\
&\quad+\bigl(y_1\omega_{13}(\alpha')+y_2\omega_{23}(\alpha')\bigr)e_3.
\end{align*}
前两项正是 $\nabla_{\alpha'}Y$，最后一项是 normal 分量，故结论成立。\hfill$\square$

\par\noindent\textbf{Definition.} 如果 curve $\alpha\colon I\to M$ 的 velocity field $\alpha'$ 沿 $\alpha$ 平行，即
\[
\nabla_{\alpha'(t)}\alpha'=0,\qquad t\in I,
\]
则称 $\alpha$ 为 geodesic（测地线）。

\medskip
考虑 $\mathbb{R}^3$ 中的嵌入 surface $M$，以及以 arc length 参数化的 curve $\alpha(s)$。曲率 $k=\|\alpha''\|$，考虑 $\alpha''$ 在 $T_{\alpha(s)}M$ 上的投影。取 adapted frame $\{e_i\}$，$e_1=\alpha'$，从而
\[
\alpha''=\frac{de_1}{ds}
=\omega_{12}(\alpha')e_2+\omega_{13}(\alpha')e_3.
\]
所以在 tangent plane 上的投影为 $\omega_{12}(\alpha')e_2$。

\textbf{Definition.} 一般地，设 $M$ 是一个 oriented 的 $2$ 维 Riemannian manifold，$\alpha\colon I\to M$ 以 arc length $s$ 参数化。沿 $\alpha$ 选取与 $M$ 的 orientation 一致的 frame $\{e_1=\alpha',e_2\}$，令
\[
k_g=(\alpha^*\omega_{12})\left(\frac{d}{ds}\right)
=\omega_{12}(\alpha'),
\]
称为 $\alpha$ 在 $M$ 中的 geodesic curvature（测地曲率）。

\medskip
\textbf{Proposition.} 设 $\alpha\colon I\to M^2$ 以 arc length 参数化，取 frame $\{e_1=\alpha',e_2\}$。则
\[
e_1\text{ 沿 }\alpha\text{ 平行}
\quad\Longleftrightarrow\quad
\alpha^*\omega_{12}=0.
\]

\textbf{Proof.}
\[
e_1\text{ 沿 }\alpha\text{ 平行}
\quad\Longleftrightarrow\quad
\nabla_{\alpha'(t)}e_1=0
\quad\Longleftrightarrow\quad
\nabla_{e_1}e_1=0.
\]
而 $\nabla_xe_1=\omega_{12}(x)e_2$，故等价性成立。\hfill$\square$

这里 $e_1=\alpha'$，所以
\[
(\alpha^*\omega_{12})_s\left(\frac{\partial}{\partial s}\right)
=(\omega_{12})_{\alpha(s)}\bigl(\alpha'(s)\bigr)
=(\omega_{12})_{\alpha(s)}\bigl(e_1(\alpha(s))\bigr),
\]
即 $\alpha^*\omega_{12}=\bigl(\omega_{12}(e_1)\circ\alpha\bigr)\,ds$。

\medskip
\textbf{Corollary.} 以 arc length 参数化的 curve $\alpha$ 是 geodesic，当且仅当其 geodesic curvature 恒为零。

\medskip
\textbf{Proposition.} 设 $M^2$ 是 oriented 的 Riemannian manifold，$\alpha\colon I\to M$ 是以 arc length $s$ 参数化的 differentiable curve，$\alpha'(s)\neq0$。设 $V$ 是沿 $\alpha$ 的非零 parallel vector field，$\varphi=\operatorname{ang}(V,\alpha')$ 为按给定 orientation 测量、沿曲线连续选取的有向角。则
\[
k_g(s)=\frac{d\varphi}{ds}.
\]

\textbf{Proof.} 令 $\bar e_1=\alpha'$，补成与 $M$ 同定向的 frame $\{\bar e_1,\bar e_2\}$；令 $e_1=V/\|V\|$，补成同定向的 frame $\{e_1,e_2\}$。局部地可将两套 frames 扩展到 $M$ 上的一个 neighborhood，并写
\[
\bar e_1=\cos\varphi\,e_1+\sin\varphi\,e_2,
\qquad
\bar e_2=-\sin\varphi\,e_1+\cos\varphi\,e_2.
\]
由 connection forms 的变换公式，
\[
\omega_{12}=\bar\omega_{12}-d\varphi.
\]
由于 $V$ 平行，$\|V\|$ 恒定，故 $e_1$ 也沿 $\alpha$ 平行。于是
\[
0=\nabla_{\alpha'}e_1=\omega_{12}(\alpha')e_2,
\qquad\text{即}\quad\omega_{12}(\alpha')=0.
\]
因此
\begin{align*}
k_g(s)
&=\bar\omega_{12}(\bar e_1)
=\omega_{12}(\alpha')+d\varphi(\alpha')\\
&=d\varphi(\alpha')
=\frac{d}{ds}(\varphi\circ\alpha)
=\frac{d\varphi}{ds}.
\end{align*}
最后使用了 chain rule；其中局部 neighborhood 上的 angle function 与沿 $\alpha$ 的 angle function 仍用同一个记号 $\varphi$。结论沿整条参数区间成立。\hfill$\square$
\end{document}














































